\section{Weighted Hilbert bounds from selected pairs}
\label{sec:hilbert-bounds}

The affine counting theorem requires lower bounds for spaces of
polynomial functions on every subvariety encountered during its
successive cuts.  We obtain these bounds by retaining a small number
of conjugate coordinate pairs.  Their image has an explicit
complete-intersection description, including at weighted infinity.

Give every \(a_i\) weight \(U>0\) and every \(b_i\) weight \(V>0\),
where \(U,V\) are fixed rational numbers.  In the transformed integral
coordinates of Proposition~\ref{prop:lattice-boxes}, these are the
weights of the entire \(x'\) and \(y'\) blocks.  Invertible linear
changes within either block preserve the weighted polynomial
filtration, even if their coefficients vary.

Let \(Y\) be an irreducible subvariety of the affine equations
\[
 a_i^kb_i=n-\theta_i\qquad(1\leq i\leq d),
\]
of dimension \(2\leq h\leq d\), with a nonempty open where all \(a_i\ne0\).
Assume that \(n\) is nonconstant.  On this open its function field
is generated by the \(a_i\) and \(n\).  Choose \(h-1\) of the
\(a_i\) algebraically independent over \(\C(n)\), and one more
distinct index.  The image in these \(h\) coordinates and \(n\)
has dimension \(h\), so its ideal is generated by one irreducible
polynomial.  Relabeling the selected indices, write it as
\[
 F_1(a_1,\ldots,a_h,n)=\sum_I a^If_I(n).
\]
None of the \(a_i\) divides \(F_1\).  Let \(Y_{\mathrm{pair}}\)
denote the closure of the projection onto the selected pairs
\((a_i,b_i)_{i=1}^h\).  It has the same function field as the
image just described.

\begin{proposition}\label{prop:weighted-hilbert}
For each nonzero coefficient polynomial \(f_I\), define
\[
 \nu_{i,I}=\operatorname{ord}_{n-\theta_i}f_I,\qquad
 e_i=\min_I(I_i+k\nu_{i,I}),
\]
and put
\begin{equation}
 D=\max_I\{U|I|+(kU+V)\deg f_I\}-U\sum_i e_i.
 \label{eq:extracted-weight}
\end{equation}
Then \(D>0\), and the weighted filtered Hilbert function of the
pair image satisfies
\begin{equation}
 H_{Y_{\mathrm{pair}}}(T)=
 \frac{(kU+V)^{h-1}D}{U^hV^h}\frac{T^h}{h!}
 +O(T^{h-1}+1).
 \label{eq:pair-hilbert}
\end{equation}
This is a lower bound for \(H_Y(T)\).
If the function field degree of \(Y\) over its pair image is at least
two, the leading coefficient in this lower bound can be doubled.
For fixed weights and bounded degree of \(Y\), all error constants
are uniform, independently of the coefficients of its equations.
\end{proposition}

\subsection{Extracting a polynomial equation}

Factor
\[
 f_I(n)=\prod_{i=1}^h(n-\theta_i)^{\nu_{i,I}}g_I(n).
\]
As in \eqref{eq:saturated-bihomogeneous-cut}, define
\begin{equation}
 C_1(a,b)=
 \sum_I a^{I+k\nu_I-e}b^{\nu_I}
 g_I(a_1^kb_1+\theta_1).
 \label{eq:affine-extracted-cut}
\end{equation}
All exponents of \(a\) are nonnegative.  Impose, in addition,
the \(h-1\) equations
\begin{equation}
 a_i^kb_i+\theta_i=a_1^kb_1+\theta_1
 \qquad(2\leq i\leq h).
 \label{eq:pair-product-relations}
\end{equation}
On the open where all \(a_i\ne0\), these equations and \(C_1=0\)
have coordinate ring
\[
 \C[a_1^{\pm1},\ldots,a_h^{\pm1},n]/(F_1),
\]
which is integral of dimension \(h\).

The boundary has smaller dimension.  If \(a_i=0\), then
\(n=\theta_i\), every other \(a_j\ne0\), and
\(b_j=(\theta_i-\theta_j)/a_j^k\).
Before imposing \(C_1\), the free coordinates are \(b_i\) and the
\(h-1\) nonzero \(a_j\), giving dimension \(h\).
The terms of \(C_1\) surviving there have
\(I_i+k\nu_{i,I}=e_i\) and are nonzero scalar multiples of
\[
 a_{-i}^{I_{-i}-e_{-i}}b_i^{\nu_{i,I}}.
\]
Their exponents determine \(I\), so these Laurent monomials are
distinct.  Thus \(C_1\) is nonzero on this stratum and reduces its
dimension to at most \(h-1\).  This checks every affine boundary
stratum, since two different \(a_i\) cannot vanish when the
\(\theta_i\) are distinct.

Every term in \eqref{eq:affine-extracted-cut} has weight at most
the number \(D\) in \eqref{eq:extracted-weight}.  The bound is
attained without cancellation.  To see this, retain the highest
weighted parts in \eqref{eq:pair-product-relations}; they say
that the products \(a_i^kb_i\) all equal a common variable \(q\).
On their nonzero-\(a\) open, the weight-\(D\) part of \(C_1\)
restricts to
\begin{equation}
 a^{-e}\operatorname{in}_{(U,kU+V)}F_1(a,q).
 \label{eq:nonzero-weighted-top}
\end{equation}
Here \(\operatorname{in}_{(U,kU+V)}\) means the sum of terms of
largest weight when each \(a_i\) has weight \(U\) and \(q\)
has weight \(kU+V\).  The expression is a nonzero Laurent
polynomial: distinct monomials \(a^Iq^j\) remain distinct after
division by \(a^e\).  This also follows directly from the
extracted-factor formula, since constants \(\theta_i\) disappear
in the highest-weight terms.
Consequently \(C_1\) has exact weight \(D\).
It is nonconstant, since otherwise the nonempty open hypersurface
defined by \(F_1\) would be empty.  Hence
\begin{equation}
 D\geq\min(U,V)>0.
 \label{eq:basic-extracted-weight}
\end{equation}

\subsection{The complete intersection at weighted infinity}

Choose a positive integer \(L_{\mathrm{wt}}\) such that
\(A=L_{\mathrm{wt}}U\) and \(B=L_{\mathrm{wt}}V\) are integers, and set
\[
 Q=kA+B,\qquad E=L_{\mathrm{wt}}D.
\]
Homogenize \eqref{eq:pair-product-relations} and \(C_1\) with a
variable \(t\) of weight one, giving \(h\) homogeneous equations
in \(2h+1\) variables.  On \(t\ne0\), scaling to \(t=1\)
identifies the scheme with the affine one times \(\C^\times\).
Its nonzero-\(a\) open is integral and generically reduced of
dimension \(h+1\).  Its remaining affine strata have dimension
at most \(h\), by the preceding boundary calculation.

At \(t=0\), the products \(a_i^kb_i\) must coincide.
If all \(a_i\ne0\), the variables \(a_1,\ldots,a_h,q\) are free
before the last equation, and \eqref{eq:nonzero-weighted-top}
is nonzero.  This part of infinity has dimension at most \(h\).
If some \(a_i=0\), the common product is zero, so
\(a_j^kb_j=0\) for every \(j\).  This locus is a union of
coordinate \(h\)-planes.  Its dimension is \(h\), even if the
last equation vanishes on an entire such plane.

It follows that the homogeneous system has dimension \(h+1\),
and no top-dimensional component is contained in either boundary.
The complete-intersection criterion used in Lemma~\ref{lem:bicone}
therefore makes its ideal prime.  In particular, it is saturated
with respect to \(t\).  Dehomogenization gives precisely the prime
ideal of \(Y_{\mathrm{pair}}\): on the nonzero-\(a\) open the
descriptions agree, and the boundary has smaller dimension.

The degree-\(N\) part of this graded quotient is the space of
affine functions represented by monomials of weight at most \(N\)
in the scaled weights \(A,B\).  Surjectivity follows by
homogenizing each such monomial with a power of \(t\).
For injectivity, a homogeneous element vanishing after
dehomogenization vanishes after localization at \(t\), and
saturation implies that it was zero in the graded quotient.
Thus the Hilbert series of these filtered dimensions is
\begin{equation}
 \frac{(1-z^Q)^{h-1}(1-z^E)}
 {(1-z)(1-z^A)^h(1-z^B)^h}.
 \label{eq:weighted-hilbert-series}
\end{equation}

At \(z=1\) this series has a pole of order \(h+1\), with leading
factor \(Q^{h-1}E/(A^hB^h)\).
We also check all other roots of unity.  If only one of
\(1-z^A,1-z^B\) vanishes, its denominator order is \(h\).
If both vanish at a nontrivial root \(\zeta\), its order divides
\(\gcd(A,B)\), and therefore divides \(Q\) and \(E\): the latter
is a monomial weight of the nonzero polynomial \(C_1\).
All \(h\) numerator factors then vanish, reducing the denominator
order \(2h\) to at most \(h\).  The extra factor \(1-z\) is
nonzero at \(\zeta\).
Partial fractions consequently give
\[
 \frac{Q^{h-1}E}{A^hB^h}\frac{N^h}{h!}
 +O(N^{h-1}+1)
\]
for the coefficient of \(z^N\).  Taking \(N=\lfloor L_{\mathrm{wt}}T\rfloor\)
and rescaling proves \eqref{eq:pair-hilbert}.
In particular, lower-dimensional coordinate planes at infinity
do not change its leading coefficient or introduce a periodic
term of order \(T^h\).

\subsection{Uniformity and function-field extensions}

The degree of \(F_1\) is bounded in terms of \(d,k\) and the degree
of \(Y\).  Indeed, \(n=a_1^kb_1+\theta_1\) is a polynomial of
fixed degree.  Form the graph of the map to the selected \(a_i\)
and \(n\), then project linearly; the usual degree bounds for
graphs and projections give the assertion.
They do not involve coefficient heights.  For fixed \(U,V\)
and a bound on that degree, only finitely many numerator degrees
\(E\) can occur in \eqref{eq:weighted-hilbert-series}.  Its
coefficient errors are therefore uniform.  This proves the
required uniformity even for moving equations and moving
blockwise linear changes of coordinates.

Pullback from the pair image is injective and preserves weights,
so its Hilbert function is a lower bound for that of \(Y\).
Let \(K_{\mathrm{pair}}=\C(Y_{\mathrm{pair}})\).
If \([\C(Y):K_{\mathrm{pair}}]\geq2\), one of the original \(a\) or \(b\)
coordinates, say \(\xi\), lies outside \(K_{\mathrm{pair}}\), since these
coordinates generate \(\C(Y)\).  Its weight \(\omega\) is
\(U\) or \(V\).  The two spaces consisting of pair-image
functions of weight at most \(T\), and \(\xi\) times such
functions of weight at most \(T-\omega\), are linearly
independent on \(Y\): a relation \(A+\xi B=0\) with
\(A,B\in K_{\mathrm{pair}}\) would contradict \(\xi\notin K_{\mathrm{pair}}\) unless
\(A=B=0\).  Hence
\begin{equation}
 H_Y(T)\geq
 H_{Y_{\mathrm{pair}}}(T)+H_{Y_{\mathrm{pair}}}(T-\omega)
 \label{eq:field-degree-gain}
\end{equation}
for \(T\geq\omega\).  The bounded shift leaves the leading
coefficient unchanged and only changes the uniform
\(O(T^{h-1}+1)\) term.  This proves the asserted doubling
and completes Proposition~\ref{prop:weighted-hilbert}.

\subsection{The general dimensional thresholds}

We record the consequences needed before the special quartic
surface argument.  Here \(Y\) is an irreducible subvariety of one of
the initial affine sets obtained in
Proposition~\ref{prop:archimedean-cuts}.  If \(Y\) has dimension \(d\), it must be
one of the mixed affine components of
Proposition~\ref{prop:archimedean-cuts}; the second-cut sets have
smaller dimension.  In this case the equation is
\(F(a,n,1)\), with \(F\) of bidegree \((l,s_0)\) and saturated
bidegree \((\lambda,s_0)\), where \(\lambda,s_0\geq1\).
Homogeneity in \(a\), and the fact that its degree in \(n\)
is exactly \(s_0\), give
\[
 D=lU+(kU+V)s_0-U\sum_i e_i
   =\lambda U+s_0V\geq U+V.
\]
Thus \eqref{eq:pair-hilbert} gives a lower bound with leading
coefficient
\[
 \frac{(kU+V)^{d-1}(U+V)}{U^dV^d}.
\]
In every dimension \(2\leq h\leq d-1\), the general bound
\eqref{eq:basic-extracted-weight} gives the coefficient
\[
 \frac{(kU+V)^{h-1}\min(U,V)}{U^hV^h}.
\]

If instead we can select \(h\) algebraically independent
\(a\)-coordinates, the same pair-image construction applies.
Its polynomial \(C_1\) must involve a \(b\)-coordinate; otherwise
it is a nonzero relation among those independent \(a\)'s.
Consequently \(D\geq V\), giving coefficient
\[
 \frac{(kU+V)^{h-1}V}{U^hV^h}.
\]
After the quartic infinite-fiber removal in
Lemma~\ref{lem:quartic-fibers}, \(n\) is algebraic over the
\(a\)-coordinate field on every remaining subvariety.
That field has transcendence degree \(h\), so the independent
selection exists.  In particular, this last coefficient supplies
the quartic \(h=3\) bound.  The next section sharpens the
remaining quartic surface case.
